\section{Síntesis y ampliación}

\subsection{Repaso de los puntos centrales}

\begin{enumerate}
    \item Los \textbf{pasos intermedios} son el mecanismo central para mejorar la capacidad de razonamiento de los LLM, ya sea mediante training, fine-tuning o prompting.
    \item \textbf{Self-Consistency} mejora de manera significativa la exactitud del razonamiento desde una perspectiva probabilística, mediante muestreos múltiples y votación mayoritaria.
    \item Las \textbf{estrategias de razonamiento} van desde few-shot CoT hasta zero-shot CoT, y de ahí a analogical reasoning y decodificación sin prompt, con formas de activación cada vez más flexibles.
    \item Las \textbf{limitaciones} —interferencia del contexto, fallas de autocorrección, sensibilidad al orden de las premisas— señalan las direcciones de mejora futuras.
\end{enumerate}

\subsection{Implicaciones para los LLM Agent}

\begin{itemize}
    \item \textbf{Llamadas a herramientas y retroalimentación externa}: un Agent puede proporcionar un Oracle externo llamando a ejecutores de código, motores de búsqueda y otras herramientas, lo cual compensa la incapacidad de los LLM para autocorregirse.
    \item \textbf{La generación aumentada por recuperación requiere cautela}: los documentos que introduce el RAG pueden contener información irrelevante que, en cambio, interfiere con el razonamiento; se necesita una recuperación y un filtrado precisos.
    \item \textbf{Organización de la información}: el orden en que se presenta la información en el flujo de trabajo de un Agent debe alinearse con la lógica del razonamiento.
    \item \textbf{Escalamiento del razonamiento}: Self-Consistency y Analogical Reasoning pueden servir como medios de escalamiento del cómputo en tiempo de inferencia.
\end{itemize}

\subsection{Lecturas adicionales}

\begin{itemize}
    \item Wei et al., ``Chain-of-Thought Prompting Elicits Reasoning in Large Language Models,'' NeurIPS 2022.
    \item Wang et al., ``Self-Consistency Improves Chain of Thought Reasoning in Language Models,'' ICLR 2023.
    \item Zhou et al., ``Least-to-Most Prompting Enables Complex Reasoning in Large Language Models,'' ICLR 2023.
    \item Kojima et al., ``Large Language Models are Zero-Shot Reasoners,'' NeurIPS 2022.
    \item Yasunaga et al., ``Large Language Models as Analogical Reasoners,'' ICLR 2024.
    \item Huang et al., ``Large Language Models Cannot Self-Correct Reasoning Yet,'' ICLR 2024.
    \item P\'{o}lya, ``How to Solve It,'' Princeton University Press, 1945.
\end{itemize}

%% --- Fin del contenido --- %%

\end{document}
